\documentclass[preprint]{glimmer-paper}

\title{Proactive Provenance: Making the Research-Object Graph Executable\\ for AI-Native Reproducibility}
\author{Shady El Damaty}
\affil{Holonym / Opscientia}
\version{0.6.0}
\repo{https://github.com/hebbianloop/glimmer}
\date{2026}

\newcommand{\code}[1]{\texttt{#1}}

\begin{document}
\maketitle

\begin{abstract}
Reproducible-pipeline tooling (BIDS, DataLad, Nipype, containers) made the \emph{bytes} of
a computational analysis recoverable, and provenance standards (W3C PROV, RO-Crate, CWL)
made its \emph{structure} describable. Yet ``this result reproduces'' remains a claim a
reader must take on trust: the description of a run and the act of re-running it are
separate artifacts, and nothing in the typed record forces them to agree. The gap widens
under AI-native workflows, where an autonomous agent re-derives the same numbers
inconsistently, repeats settled data-processing errors, and loses prior findings across
sessions --- each a reproducibility failure, because the agent has no executable contract
to anchor to. We introduce \emph{proactive provenance}: a typed research-object graph whose
runs are first-class, \textbf{executable, standard-gated, and self-verifying} nodes. We add
to Glimmer --- a per-entity-sidecar research-object knowledge base --- a single core node
type, the \code{run-record} (one concrete, replayable invocation, in the sense of a PROV
\emph{Activity}), and a node runner (\code{glimmer run} / \code{glimmer rerun}) that, for
each run-record: gates its inputs (they must be both content-pinned \emph{and} valid against
their declared standards, with validation delegated to standard-specific validators),
replays the recorded command in a pinned container, and verifies the outputs at one of three
fidelity tiers --- \emph{byte-identical} (with header normalization), \emph{numeric-within-tolerance}
(re-deriving a reported number from source), and \emph{structural} (for agent/LLM outputs).
The run-record is the executable unit of the agentic research loop: a hypothesis decomposes
into planned runs, running them emits derivatives and findings with recorded verdicts, and
those verdicts --- not memory --- are what the next planning step reads. We describe the
design, a reference implementation, and a worked example exercising all three tiers, the
standards gate, and runtime-certified method equivalence with no specialized dependencies.
\end{abstract}

\section{Introduction}

The 2010s solved a narrow version of reproducibility: with versioned data
(DataLad / git-annex \citep{halchenko2021datalad}), versioned code, and pinned environments
(containers), the \emph{bytes} an analysis consumed and produced can be recovered later. The
2020s added a \emph{description} layer --- W3C PROV activities \citep{provdm2013}, RO-Crate
manifests \citep{soilandreyes2022rocrate}, CWL tool/workflow descriptions
\citep{crusoe2022cwl} --- so the \emph{structure} of a computation is machine-readable.
Glimmer's prior contribution \citep{eldamaty2026glimmer} sat above both: a typed-entity
graph (datasets, methods, derivatives, findings, standards, publications, concepts)
distributed across per-entity sidecars, so an agent could \emph{navigate} provenance and
render auditable decisions.

But all of this is \textbf{descriptive}. A \code{derivative} node records an
\code{output-hash}; a \code{method} records a tool version and a workflow SHA; PROV records
that activity $A$ used entity $E$ and generated entity $O$. None of it re-runs $A$ and checks
that $O$ still results. The description of the run and the run itself are different objects,
maintained by hand, and free to drift.

AI-native workflows make the gap acute. An autonomous or semi-autonomous agent operating
over a dataset, with no executable contract to anchor to, exhibits a characteristic failure
pattern: it re-derives the same quantity by different routes and gets different answers,
repeats data-processing mistakes that were already diagnosed, and forgets findings
established earlier in the same project. Each is a reproducibility failure, and each is a
direct consequence of the substrate being a \emph{ledger of claims} rather than a \emph{set
of runnable, self-checking contracts}.

\paragraph{Thesis.} Reproducibility should be an \emph{executable property of the graph}, not
a property of a separate pipeline a reader is invited to trust. We call this \textbf{proactive
provenance}: the research-object graph not only records what an output is, but can re-run the
act that produced it and verify the result --- at a fidelity appropriate to the
computation's nature. We realize it with one new node type and one tool, and argue this is
the missing primitive that makes the agentic research loop self-sustaining.

\paragraph{Contributions.}
\begin{enumerate}
\item The \code{run-record} node type: a first-class, lifecycle-bearing PROV \emph{Activity}
  that binds a method, pinned and standard-validated inputs, expected outputs, the exact
  command, and a pinned environment --- the executable unit of the agentic loop
  (\S\ref{sec:runrecord}, \S\ref{sec:loop}).
\item A node runner with a pre-run \textbf{standards gate} and \textbf{three verification
  tiers}, including a faithful byte-identical tier that normalizes non-semantic header bytes
  (\S\ref{sec:runrecord}).
\item \textbf{Standards as a runtime gate}: the run-record \emph{references} the standards
  its inputs must satisfy; the runner \emph{enforces}; validators are \emph{delegated}
  (\S\ref{sec:gate}).
\item \textbf{Runtime-certified method equivalence}: two implementations are
  \code{equivalent-to} only if the runner confirms they produce matching outputs on shared
  inputs (\S\ref{sec:registry}).
\item A reference implementation and a dependency-light worked example demonstrating all of
  the above (\S\ref{sec:impl}).
\end{enumerate}

\section{Background and related work}

\paragraph{Reproducible pipelines.} BIDS \citep{gorgolewski2016bids} standardizes
neuroimaging data layout; DataLad/git-annex \citep{halchenko2021datalad} version data and
pin content by hash; Nipype \citep{gorgolewski2011nipype} and container images (e.g.
fMRIPrep \citep{esteban2019fmriprep}) pin the toolchain. These guarantee \emph{re-fetch} and
\emph{re-execution capability} but not \emph{verification}: nothing asserts that
re-execution reproduces the recorded result. Glimmer reuses this layer wholesale --- DataLad
coordinates live on every node --- and adds the act of checking.

\paragraph{Provenance description.} W3C PROV \citep{provdm2013} models
entities/activities/agents; RO-Crate \citep{soilandreyes2022rocrate} packages a dataset with
typed metadata; CWL and \code{cwlprov} \citep{crusoe2022cwl} describe and record runs. The
\code{run-record} is deliberately PROV-shaped (an Activity that \code{consumes} entities and
\code{regenerates} derivatives), but it is \emph{executable in place}: the same node a reader
inspects is the node the runner replays and stamps with a verdict. The contribution is not a
new description format but closing the loop between description and execution inside one
typed object.

\paragraph{Workflow re-execution.} \code{datalad run} / \code{datalad containers-run} record
a command with its inputs/outputs and can re-run it; CI harnesses re-execute pipelines. The
node runner builds directly on \code{datalad containers-run} for replay. What Glimmer adds is
(a)~re-execution as a \emph{graph operation} over typed nodes rather than a shell convention,
(b)~a \emph{standards gate} before execution, and (c)~\emph{tiered} verification spanning
deterministic, stochastic, and agent-produced outputs --- not only byte-equality.

\paragraph{AI-for-science systems.} End-to-end autoresearch systems
\citep{lu2024aiscientist} run fixed internal pipelines. Glimmer is complementary substrate:
the run-record gives such systems a typed, replayable, verifiable unit of work that survives
independently of any one system, so multiple agents can cooperate on, and audit, the same
graph.

\section{The Glimmer graph (recap)}

Glimmer represents a research project as a graph of typed nodes, each a Markdown file with
YAML front-matter (a ``sidecar''), enumerated by a root index. Core node types include
\code{dataset}, \code{method}, \code{derivative}, \code{finding}, \code{standard},
\code{publication}, \code{concept}, \code{experiment}, \code{instrument}, \code{persona},
\code{organization}, and \code{program}; edges are properties on the source node.
Domain-specific vocabulary lives in \emph{profiles} (e.g. neuroimaging/BIDS) that augment
core types without forking the core. Every node carries a content hash and its DataLad
re-fetch coordinates; an agent that produces a node must attach a \code{reasoning-trace}
citing the nodes it read. The graph is plain files, so it survives \code{git clone},
\code{datalad export}, and \code{rsync} with no bespoke database. This much was descriptive;
the next section makes it executable.

\section{Proactive provenance: the \texttt{run-record}}
\label{sec:runrecord}

\subsection{The node}
A \code{method} is a \emph{reusable} tool (one skull-strip, applied to many subjects); a
\code{derivative} is a \emph{product} (a file and its hash). Neither captures \emph{the act}.
The \code{run-record} is that act --- one concrete invocation: \textbf{this} command, on
\textbf{these} content-pinned inputs, in \textbf{this} container, testing \textbf{this}
hypothesis, on \textbf{this} date. Its required core is a \code{method}, a \code{command}, a
\code{provenance-mode} (which selects the default verification tier), a \code{status}, a list
of pinned \code{inputs}, and a list of expected \code{outputs}; edges (\code{reruns},
\code{consumes}, \code{regenerates}, \code{emits}, plus \code{tests-hypothesis} /
\code{addresses-concept}) make it navigable.

It carries a \textbf{lifecycle}: \code{planned} (written when a hypothesis is decomposed;
inputs may be a \emph{specification} rather than a pin) $\rightarrow$ \code{ready} (inputs
pinned and standard-valid: the gate passed) $\rightarrow$ \code{running} $\rightarrow$
\code{executed} (outputs hashed, verdict written) / \code{failed} / \code{superseded}. The
runner advances the lifecycle and writes two blocks back into the node: \code{validation-gate}
(the pre-run check) and \code{replay-verdict} (the per-output result).

\subsection{The node runner}
One engine, two modes. \code{glimmer run} is \emph{forward execution} --- gate, replay, hash
and record outputs, advance \code{ready}~$\rightarrow$~\code{executed}. \code{glimmer rerun}
is \emph{reproduction} --- re-execute an already-recorded run and compare its outputs to what
was recorded. The runner is for running; reproduction is the special case of re-running
something already recorded. Per run-record the engine (i)~resolves and pins inputs
(materializing via \code{datalad get}, matching the recorded annex-key/commit-sha; an
unreachable input short-circuits to an \code{inputs-unavailable} verdict that records data
liveness --- never a false pass), (ii)~enforces the standards gate (\S\ref{sec:gate}),
(iii)~replays the command pinned to the container digest via \code{datalad containers-run},
falling back to \code{datalad run} and then a host subprocess flagged as a dirty environment,
(iv)~verifies outputs per tier (\S\ref{sec:tiers}), and (v)~emits a provenance manifest and,
optionally, writes the verdict back into the graph.

\subsection{Standards as a runtime gate}
\label{sec:gate}
A run-record only \textbf{references} the standards its inputs must satisfy (an input-pin's
\code{conforms-to}, the \code{requires-standard} edge). The runner \textbf{enforces} that
each input is both \emph{pinned} and \emph{valid} before any command runs; a failure yields a
\code{gate-failed} verdict and no execution. The check itself is \textbf{delegated} to the
standard's \code{validator} hint (e.g. \code{bids-validator}), or falls back to graph-level
validation, recording \code{unchecked} honestly when no validator is available. The
run-record thus never reimplements validation, yet ``ran on standard-valid data'' becomes an
enforced, recorded fact --- moving standards from passive metadata into the runtime.

\subsection{Three verification tiers}
\label{sec:tiers}
Reproducibility is not one thing, and a single notion of ``match'' is wrong for most science.
The tier is selected by \code{provenance-mode} and overridable per output.
\begin{itemize}
\item \textbf{Byte-identical} (deterministic). Re-execute; the output hash must match exactly
  --- \emph{after normalization}, because two correct re-runs of an FSL/Nipype step differ in
  volatile header bytes (timestamps, \code{descrip} fields) that carry no scientific content.
  The runner zeroes those fields for NIfTI (keeping affine, dtype, and voxels exactly),
  strips GIFTI metadata, and canonicalizes JSON, recording which normalization was applied;
  if the imaging library is unavailable it falls back to a raw compare, recorded as such ---
  never a silent pass.
\item \textbf{Numeric-within-tolerance} (stochastic). Re-derive a number from source and
  assert $|\,\text{observed}-\text{expected}\,| \le \max(\text{abs},\,\text{rel}\cdot|\text{expected}|)$.
  This is the honest guarantee for stochastic analyses, and the mechanism for
  \textbf{reproducing a paper's reported numbers from source} within a declared tolerance.
\item \textbf{Structural} (agent-inferred). LLM/agent prose is not hash-reproducible.
  Verification is structural: the node's \code{reasoning-trace} must cite real graph nodes,
  those nodes must contain the values the trace claims, and any retriever manifest must be
  self-consistent --- the analogue of byte-equality for the inferential regime.
\end{itemize}
Each verdict is recorded with its tier and degradations, so a ``verified'' verdict is never
silently weaker than its claim (a host-fallback run is marked as such; a tampered expected
hash yields \code{mismatch}).

\subsection{Method equivalence and a cross-project registry}
\label{sec:registry}
A method is a reusable function over pinned data, and its \emph{semantic pattern} --- ``compute
the mean of a series'', ``skull-strip a T1w'' --- is not unique to one project. We separate
pattern from implementation with a \code{registry-ref} (a cross-project namespaced pattern
id) and the edges \code{implements}, \code{equivalent-to}, and \code{refines}. Crucially,
\code{equivalent-to} is \textbf{certified, not asserted}: the runner runs both
implementations on the same pinned inputs and confirms their outputs match (byte-identically
or within tolerance). This makes ``an equivalent, cleaner, or faster implementation is
acceptable'' a checkable statement, and lets a method's identity --- and the verification
baseline attached to it --- outlive a single graph. A dedicated abstract
\code{method-pattern} node type and a published cross-institution registry are future work
(\S\ref{sec:future}).

\section{The agentic loop, made executable}
\label{sec:loop}
Glimmer previously specified a plan~$\rightarrow$~run~$\rightarrow$~feedback~$\rightarrow$~replan
research loop --- a \code{concept} decomposed into hypotheses, agents producing derivatives
and findings, a human reviewing --- but it had no runtime primitive; the ``run'' step left
only a derivative behind. The run-record \emph{is} that primitive. A hypothesis acquires one
or more \code{planned} run-records (\code{tests-hypothesis}~$\rightarrow$~\code{concept});
running them gates, executes, and emits \code{derivative}s (\code{regenerates}) and a
\code{finding} (\code{emits}) whose \code{addresses-concept} edge closes the loop back onto
the question. The feedback the next iteration reads is the recorded verdict plus the emitted
finding --- not the agent's memory. This is precisely why the loop becomes self-sustaining and
why the agent failure modes of \S1 are structurally suppressed: a finding cannot be silently
lost (it is a node addressing the concept), and a settled mistake cannot be silently repeated
(the prior run's verdict is in the graph and a re-run must reproduce it).

\section{Implementation and worked example}
\label{sec:impl}
The implementation is roughly 600 lines of dependency-light Python (PyYAML only; DataLad, a
container runtime, the imaging library, and external validators are all feature-detected and
degrade with the degradation recorded). The schema change is additive and
backward-compatible; existing graphs validate unchanged.

The worked example (\code{examples/synthetic-provenance/}) is the smallest artifact that
exercises the whole design with no real data and no specialized dependencies. A \code{concept}
is decomposed into four planned runs over a fixed synthetic signal whose input declares
conformance to a toy standard with an executable validator (Table~\ref{tab:results}).

\begin{table}[h]
\centering
\small
\begin{tabular}{@{}llll@{}}
\toprule
run & provenance-mode & tier & reproduce verdict \\
\midrule
\code{run-synth-mean}          & deterministic   & byte-identical          & \code{verified} \\
\code{run-synth-mean-fast}     & deterministic   & byte-identical          & \code{verified} (certified equiv.) \\
\code{run-synth-classifier}    & stochastic      & numeric ($\pm 0.02$)    & \code{reproduced-within-tolerance} \\
\code{run-synth-agent-summary} & agent-inferred  & structural              & \code{structurally-valid} \\
\bottomrule
\end{tabular}
\caption{The worked example exercises all three tiers; reproduction is 100\%.}
\label{tab:results}
\end{table}

Forward \code{glimmer run} advances the planned runs to executed and emits a finding that
addresses the concept; \code{glimmer rerun} reproduces all four at 100\%. Negative controls
behave as designed: a malformed input is caught by the standards gate before execution
(\code{gate-failed}); a tampered expected hash yields \code{mismatch}; an unreachable input
under \code{-{}-offline} yields \code{inputs-unavailable}; bypassing the gate is recorded as a
dirty run; and equivalence certification confirms the two mean implementations produce
identical output. Each maps to a non-zero exit, so the example doubles as a CI check.

\section{Discussion}
Proactive provenance reframes reproducibility from a property a reader audits by hand to a
property the substrate enforces and re-checks. Three design choices carry the weight.
\emph{First}, tiered verification: insisting on byte-equality for everything is both too
strong (stochastic analyses fail it) and too weak (it ignores agent outputs entirely);
matching the check to the computation's nature is what lets one mechanism span a deterministic
pipeline, a reported statistic, and an LLM interpretation. \emph{Second}, delegated gating:
standards become runtime-enforceable without the runner knowing any standard's internals, so
the design rides the existing validator ecosystem rather than reimplementing it.
\emph{Third}, certified equivalence: it converts the social convention ``use an equivalent
method'' into a machine-checked relation, the seed of a shared verification baseline across
institutions. Together these make the claim ``AI-native reproducibility'' concrete: an
agent's unit of work is now a contract it can execute and that others can re-execute, with a
verdict legible to both machines and reviewers.

\section{Limitations and future work}
\label{sec:future}
The runner replays a recorded command; it does not yet \emph{generate} a reproduction path
for an arbitrary external paper. The planned-run-record with specified (unpinned, possibly
surrogate) inputs is the affordance for \textbf{minimal-path reproduction} --- synthesizing a
minimal runnable graph to test a published claim even when the original data is unavailable
--- the principal next step. The method registry currently ships as lightweight edges plus
runtime-certified equivalence; a dedicated abstract \code{method-pattern} type and a
published, cross-institution registry (with a reputation/provenance model for who contributed
which pattern or baseline) are deferred. Container replay requires DataLad and a runtime; on a
bare host the runner degrades and records the weaker environment, but does not reconstruct it.
Structural verification establishes that an agent cited real values, not that its
interpretation is correct --- it guarantees auditability, not truth. Finally, the worked
example is synthetic by design; validating the byte-identical and numeric tiers on a full
neuroimaging pipeline and on a dissertation's reported numbers is in progress in a downstream
project that imports the runner as a dependency.

\section{Conclusion}
We argued that reproducibility belongs \emph{inside} the typed research object as an
executable property, and realized it with a single node type and one tool. The
\code{run-record} makes a run a first-class, lifecycle-bearing, standard-gated, self-verifying
graph node; the node runner enforces the gate and verifies outputs at a fidelity matched to
the computation. The result turns the agentic research loop from a description into a running
system whose feedback is recorded verdicts rather than recollection --- a concrete substrate
for AI-native science in which ``this reproduces'' is something the machine demonstrates, not
something the reader is asked to believe.

\bibliographystyle{plainnat}
\bibliography{refs}

\end{document}
